\chapter{Project Execution}
\label{chap:execution}

\noindent
This chapter describes the systematic implementation and execution of the
reproduction experiments for \cite{hoyer2023ijcv_sde} and the proposed
cross-task consistency loss extension. Beyond a log of execution, this
chapter delineates the rationale behind dataset configurations, the
architecture of the training pipeline, the design of the PAD multi-task
decoder, and the orchestration of computational resources on a
high-performance computing (HPC) cluster. Throughout, this study adopts
the paper's terminology: \textbf{Stage~(a)} refers to standard supervised
learning and data selection, while \textbf{Stage~(b)} encompasses
semi-supervised mechanisms and multi-task learning.

% ============================================================================
\section{Experimental Setup}
\label{sec:setup}

\subsection{Hardware and Software Environment}

This project was executed across two successive HPC environments,
reflecting a mid-project migration in available computing resources.
\begin{itemize}
    \item \textbf{Reproduction experiments (exp~210--213, Tables~1,
    3, 5, 7)}: These runs were submitted to
    Isambard-AI~\cite{isambard_ai}, the University of Bristol's
    AI-optimised supercomputing facility, using a single GPU per job.
    \item \textbf{CT loss extension (exp~219--221)}: All
    cross-task consistency loss experiments were executed on the
    University of Bristol's BluePebble~(BP1) cluster under the SLURM
    workload manager, following the migration to this newer resource.
    BP1 is a heterogeneous cluster; Table~\ref{tab:bp1_nodes}
    summarises the GPU node types encountered during these runs.
\end{itemize}

\begin{table}[H]
\centering
\caption{BP1 GPU node types allocated to CT loss experiments (exp~219--221).}
\label{tab:bp1_nodes}
\renewcommand{\arraystretch}{1.2}
\begin{tabular}{llcc}
\toprule
\textbf{Node class} & \textbf{GPU model} & \textbf{GPU VRAM} &
\textbf{Observed usage} \\
\midrule
\texttt{bp1-gpu014} & NVIDIA RTX~2080 & 8\,GB  & minority of runs \\
\texttt{bp1-gpu030} & NVIDIA RTX~3090 & 24\,GB & baseline MTL reruns \\
\texttt{bp1-gpu040} & NVIDIA V100     & 32\,GB & majority of CT loss runs \\
\bottomrule
\end{tabular}
\end{table}

\noindent
Job-to-node assignment on BP1 is managed by SLURM and is not directly
controllable; the allocation therefore varies across submissions. The
Phase~I and Phase~II CT loss runs (exp~219, 220) executed predominantly
on \texttt{bp1-gpu040}~(V100, 32\,GB VRAM), which is the most VRAM-capable
node class available on BP1 and the relevant reference hardware for
the overhead analysis in Chapter~\ref{chap:evaluation}. The
implementation was isolated within a dedicated Conda environment with
pinned dependencies to ensure consistent execution across node types.

\paragraph{Resource Allocation per Job (BP1).}
All CT loss experiments on BluePebble were allocated the following
resources per SLURM job:
\begin{itemize}
    \item \textbf{GPU}: One GPU per job (\texttt{--gres=gpu:1}); node
    type determined at dispatch time.
    \item \textbf{CPU and Memory}: 8 physical cores and 64\,GB of system
    RAM per job, providing sufficient capacity for multi-worker data
    pre-fetching without I/O bottleneck.
    \item \textbf{Wall-time}: 24\,hours for Phase~I single-seed runs;
    40\,hours for Phase~II (exp~220) and exp~221 three-seed runs, with
    the three seeds executed sequentially within a single allocation.
\end{itemize}

\paragraph{Software and Reproducibility Controls.}
Random seeds are fixed across Python, NumPy, and PyTorch at the start of
every run. The full configuration of each experiment---experiment ID,
label budget, seed, and all hyperparameters---is logged to disk alongside
TensorBoard event files and model checkpoints. The codebase is
version-controlled, enabling any run to be exactly reproduced from its
recorded experiment ID.

\paragraph{Training Protocol.}
All reproduction and extension experiments use a ResNet-101 encoder
pretrained on ImageNet. Stage~(b) and cross-task loss experiments
additionally employ the PAD multi-task decoder
(Section~\ref{sec:pad-decoder}). Training uses the SGD optimiser (momentum
$0.9$, weight decay $5\!\times\!10^{-4}$) with a polynomial
learning-rate decay schedule (\texttt{stepx}) over 40{,}000 iterations.
Learning rates are set to $10^{-2}$ for the decoder, $10^{-3}$ for the
backbone encoder, and $10^{-6}$ for the pose network. Due to the
single-GPU constraint, a batch size of 2 is used without Synchronised
Batch Normalisation~(SyncBN); this is expected to introduce minor
discrepancies relative to the multi-GPU setup of the original paper, as
each normalisation layer receives more limited batch statistics.

% ============================================================================
\section{The Dataset and Feature Representation}
\label{sec:dataset}

The pipeline exploits two complementary data modalities within the
Cityscapes dataset~\cite{cordts2016cityscapes}: pixel-level semantic
annotations for supervised segmentation and adjacent temporal frames for
geometric self-supervision.

\subsection{Data Composition}
\begin{itemize}
    \item \textbf{Fine Annotations}: 2{,}975 fine-annotated images
    constitute the labelled pool. Low-label budgets of 100, 372, and 744
    samples are simulated by selecting subsets from this pool. The
    author-provided fixed subset indices are used throughout to maintain
    exact parity with the original evaluation protocol.
    \item \textbf{Adjacent Temporal Sequences}: For each target image, a
    3-frame window (current frame~0, and neighbouring frames $-1$ and
    $+1$) supplies geometric cues for monocular self-supervised depth
    estimation and geometry-aware components such as DepthMix~(DX).
    These frames carry no semantic annotations and are used only for
    photometric and pose supervision signals.
\end{itemize}

\subsection{Geometric Constraints and Preprocessing}
Images are resized to $1024\!\times\!512$ and randomly cropped to
$512\!\times\!512$ during training. This resolution balances sufficient
spatial context for boundary-sensitive tasks against single-GPU VRAM
feasibility, particularly for Stage~(b) configurations that maintain two
parallel decoder branches simultaneously.

\subsection{SDE Checkpoints for Downstream Components (fd0/fd2)}
\label{sec:fd-checkpoints}
Several Stage~(b) components require an SDE-pretrained encoder to
provide depth-informed feature initialisations. \texttt{fd0} refers to
encoder weights from standard monocular depth pretraining; \texttt{fd2}
refers to weights from feature-distance regularised pretraining, which
additionally prevents forgetting of geometric representations during
downstream fine-tuning. To maintain a controlled reproduction setting,
we use the author-released checkpoints and keep them fixed across all
downstream experiments.

% ============================================================================
\section{The PAD Multi-Task Decoder}
\label{sec:pad-decoder}

Stage~(b) experiments and the proposed $\mathcal{L}_\text{ct}$ extension
rely on the Pyramid Aggregation and Decoupling~(PAD) decoder introduced
in \cite{hoyer2023ijcv_sde}. PAD extends a standard depth decoder with a
parallel segmentation branch that shares the ResNet-101 encoder features
but maintains task-specific decoder parameters, enabling joint
optimisation of depth estimation and semantic segmentation.

The architecture proceeds in two execution phases around an intermediate
\emph{distillation layer} (decoder upsampling layer~7 in our
configuration, producing feature maps of spatial resolution
$128\!\times\!128$ with 128 channels):
\begin{enumerate}
    \item \textbf{First half (encoder to distillation layer)}: The depth
    and segmentation decoder branches independently process the shared
    encoder features up to the distillation layer. Self-attention
    modules~(SA) are applied to each branch's intermediate feature map
    to capture long-range spatial dependencies. The attention-weighted
    depth features are then added to the segmentation branch state, and
    vice versa, enabling bidirectional cross-task context exchange before
    the second half begins.
    \item \textbf{Second half (distillation layer to output)}: The depth
    and segmentation decoders continue from the merged representations to
    their respective final prediction heads, producing depth maps and
    per-pixel segmentation logits.
\end{enumerate}

The intermediate feature maps at the distillation layer---denoted
$f_\text{seg}$ and $f_\text{depth}$---are precisely the representations
that the proposed $\mathcal{L}_\text{ct}$ acts upon
(Chapter~\ref{chap:evaluation}, Section~\ref{sec:ct-calibration-ablation}).
This placement is deliberate: the distillation layer sits after the
self-attention exchange, so both feature maps already encode cross-task
contextual information, making the explicit alignment introduced by
$\mathcal{L}_\text{ct}$ a targeted refinement of an already-informed
representation rather than a cold alignment from scratch.

% ============================================================================
\section{Configuration Architecture: The Experiments Hub}
\label{sec:design}

To manage the high dimensionality of hyperparameters across multiple
reproduction tables and ablation studies, a centralised
\texttt{experiments.py} module serves as a single source of truth,
mapping discrete experiment IDs to fully-specified hyperparameter
dictionaries. This prevents configuration drift between runs and enables
atomic reproducibility of any individual result by ID.

\begin{itemize}
    \item \textbf{exp~210}: Baseline and mixing strategy variants for
    Table~3 in \cite{hoyer2023ijcv_sde}, plus baseline and transfer
    variants for Table~5.
    \item \textbf{exp~211}: SDE-guided data selection variants for
    Table~1 in \cite{hoyer2023ijcv_sde}.
    \item \textbf{exp~212}: Multi-task learning~(MTL) configurations
    for Stage~(b) evaluation (Tables~5 and~7).
    \item \textbf{exp~213}: Full component ablation (S\,+\,DX\,+\,MTL)
    for Table~7 in \cite{hoyer2023ijcv_sde}.
    \item \textbf{exp~219}: Phase~I cross-task consistency loss sweep. A
    systematic single-seed (seed~7) scan over all 32 combinations of
    loss type (cosine~/ normalised MSE), projection layout (single~/
    dual), warmup schedule ($w0$\,/\,$w5k$), and
    $\lambda_\text{ct}\in\{0.25, 0.50, 0.75, 1.00\}$, each for 40{,}000
    iterations.
    \item \textbf{exp~220}: Phase~II three-seed validation. The top
    candidates from exp~219 are retrained with three independent seeds
    (7, 25, 42) to obtain statistically reliable mean and variance
    estimates. Dual-projection variants are included to enable a fair
    architectural comparison under multi-seed conditions.
    \item \textbf{exp~221}: Extended PAD-MTL component ablation. Three
    additional combinations of DepthMix~(DX) and label selection~(S)
    applied on top of the PAD-MTL baseline are validated with three
    independent seeds: MTL\,+\,DX; MTL\,+\,S; and
    MTL\,+\,DX\,+\,S. Results for exp~221 are pending completion of
    scheduled HPC maintenance and will be incorporated into
    Chapter~\ref{chap:evaluation} upon availability.
\end{itemize}

% ============================================================================
\section{Run Manifest: From Table~1 to Table~7 and Beyond}
\label{sec:hpc_exec}

\subsection{Table~1: Data Selection (exp~211)}
This task curates informative labelled subsets under fixed budgets using
three selection criteria: (i)~Random sampling as baseline; (ii)~Entropy
sampling, which ranks unlabelled images by model prediction uncertainty;
and (iii)~SDE-guided heuristics combining Uncertainty Sampling~(US) and
Diversity Sampling~(DS). For each criterion, images are scored and the
top-$K$ are selected to match the target label budget (100, 372, or 744
samples). The author-provided fixed subset indices ensure exact
comparisons across runs.

\subsection{Table~3: Mixing Strategy (exp~210)}
This run evaluates four training strategies under the 372-label and
full-supervision settings: (i)~supervised-only baseline; (ii)~pseudo-label
self-training; (iii)~ClassMix; and (iv)~DepthMix~(DX). Semi-supervised
variants follow the official teacher--student framework with
exponential-moving-average (EMA) teacher updates and confidence-filtered
pseudo-label generation. DepthMix replaces the binary masks of ClassMix
with geometry-consistent composition masks derived from SDE-generated
depth maps, reducing boundary artefacts at class transitions.

\subsection{Tables~5 and~7: Feature Transfer and Component Ablation
(exp~212, exp~213)}
These runs evaluate SDE-to-segmentation feature transfer using fd0 and
fd2 checkpoints~(Table~5), and the synergistic effects of combining data
selection~(S), DepthMix~(DX), and multi-task learning~(MTL)~(Table~7).
Table~7 requires eight distinct configurations to isolate additive
contributions; each is executed across three independent seeds (7, 25,
42) and results are reported as mean\,$\pm$\,std.

\subsection{Cross-Task Consistency Loss (exp~219, exp~220)}
\label{sec:ct_exec}

The cross-task consistency loss experiments are structured as a
two-phase protocol designed to maximise information gain per GPU~hour,
balancing thorough hyperparameter exploration against the constraints of
an academic single-GPU computing budget.

\paragraph{Phase~I --- Single-seed calibration sweep (exp~219).}
The hyperparameter space for $\mathcal{L}_\text{ct}$ comprises four
factors (loss type, projection layout, warmup schedule,
$\lambda_\text{ct}$) yielding 32 distinct configurations. Evaluating
all 32 with three seeds would require 96 full training runs of
$\sim$10 hours each---approximately 960~GPU hours, which is prohibitive
for a single-GPU academic project.

Phase~I therefore fixes seed~7 and trains all 32 configurations for the
full 40{,}000 iterations. Seed~7 is chosen as the median-performing seed
in the PAD-MTL baseline ($59.84$\,\%, compared to $60.92$\,\% and
$61.50$\,\% for seeds~25 and~42), providing a conservative reference
that avoids inflating apparent gains. Since every configuration is
evaluated under identical conditions---same seed, same data split, same
training schedule---the relative ranking across configurations is
informative even though absolute values carry single-seed variance.

The goal of Phase~I is not to identify the globally optimal
configuration, but to efficiently eliminate underperforming regions of
the search space. Configurations that show clear gains over the seed-7
baseline and represent distinct regions of the hyperparameter grid are
retained as candidates for Phase~II; the remainder are discarded without
spending additional GPU resources on multi-seed validation.

\paragraph{Phase~II --- Three-seed validation (exp~220).}
\label{par:phase2_rationale}
Multi-seed training is reserved exclusively for the candidates
identified in Phase~I. The rationale is as follows. The purpose of
running three independent seeds is to determine whether an observed
improvement is genuine and reproducible, or whether it is an artefact
of a particular random initialisation. This question is only
scientifically relevant for configurations that Phase~I has already
shown to be competitive: committing three seeds to configurations that
Phase~I ranked consistently below the baseline would consume three
times the compute while yielding no additional insight, since the
low Phase~I performance under a controlled single-seed comparison
already constitutes strong evidence of underperformance.

Selecting candidates based on Phase~I ranking is therefore not
cherry-picking---it is a principled allocation of a finite compute
budget. The selection ensures that Phase~II provides dense, statistically
reliable coverage precisely where it matters: the high-performing region
of the hyperparameter space, where genuine differences between
configurations are small enough that single-seed evidence alone is
insufficient to draw firm conclusions. Dual-projection variants are
included in Phase~II despite their Phase~I results being marginally
below the single-projection peaks, specifically to provide a fair
architectural comparison under multi-seed conditions rather than
dismissing them on single-seed evidence alone.

Each Phase~II configuration is submitted as an independent SLURM job
running three seeds (7, 25, 42) sequentially within a 40-hour
wall-time allocation. The PAD-MTL baseline (no $\mathcal{L}_\text{ct}$)
is retrained with the same three seeds in the same submission batch,
providing a consistent per-seed reference for paired comparisons and
eliminating confounds from cross-run baseline drift.

\subsection{Extended Component Ablation (exp~221)}
\label{sec:exp221}

Exp~221 complements the Table~7 reproduction by providing additional
three-seed validation for PAD-MTL component combinations that are
directly relevant to contextualising the $\mathcal{L}_\text{ct}$ results.
Three configurations are evaluated: MTL\,+\,DX (multi-task with
DepthMix augmentation), MTL\,+\,S (multi-task with SDE-guided label
selection), and MTL\,+\,DX\,+\,S (the full combination). Each is
trained with seeds 7, 25, and~42 for 40{,}000 iterations on the
372-label regime under the same pipeline as exp~219 and~220, ensuring
direct comparability with the CT loss results.

At the time of writing, exp~221 jobs have been submitted to the BP1
cluster and are queued pending completion of scheduled HPC maintenance.
Results will be incorporated into Chapter~\ref{chap:evaluation} upon
availability.

% ============================================================================
\section{Summary}
\label{sec:exec_summary}

This chapter described a controlled execution environment in which data
splits, model configurations, and HPC resources were consistently managed
throughout. By leveraging the University of Bristol's BluePebble cluster
and a centralised configuration hub (\texttt{experiments.py}), the core
reproduction experiments corresponding to Tables~1, 3, 5, and~7 of
\cite{hoyer2023ijcv_sde} were executed in a principled and reproducible
manner.

The cross-task consistency loss investigation followed a two-phase
protocol designed to maximise scientific return per GPU~hour: a
single-seed scan (Phase~I, exp~219) efficiently identified promising
regions of the 32-configuration hyperparameter space, and a three-seed
validation (Phase~II, exp~220) confirmed the statistical reliability of
the top candidates. The rationale for restricting multi-seed resources
to Phase~I candidates---rather than exhaustively validating all
configurations---is grounded in the principle that multi-seed validation
is only informative for configurations already demonstrated to be
competitive. The PAD multi-task decoder (Section~\ref{sec:pad-decoder})
provides the architectural foundation for both the reproduced MTL
baseline and the proposed $\mathcal{L}_\text{ct}$ extension. A further
ablation (exp~221) will contextualise the CT loss results against
additional PAD-MTL component combinations once HPC resources are
restored.
